\section{Literature Analysis}\label{sec:literature_analysis}

Our study focuses on modeling \ac{fc} rate (MT/day) of a bulk carrier using \ac{bbm} and fusing operational data with environmental factors. \\
In this regard, the most recent review papers in the space are done by Huang et al.~\citep{huangMachineLearningSustainable2022} and Fan et al.~\citep{fanReviewShipFuel}.  
Huang et al.~\citep{huangMachineLearningSustainable2022} covered applications of \ac{bbm} in different ship phases such as ship design, voyage planning, and operational performance. 
While Fan et al.~\citep{fanReviewShipFuel} analyzed the literature from 2001 to 2021 and focused on applications of \ac{bbm} in operational context only with a comparison of different 
\ac{ml} classifications in various ship operation stages and based on data availability.

To avoid duplication, this section will focus on work that emerged since then that is not covered in the above reviews regarding \ac{fc} prediction by fusing external data sources. 

Xie et al.~\citep{xieFuelConsumptionPrediction2023} compared between \ac{bbm} and \ac{wbm} and proved \ac{bbm} superiority. They developed two ensemble models, namely \ac{rf} and 
\ac{xgboost} models for an oil tanker vessel based on eight features. They fused external public weather data from \ac{ecmwf} considering only two parameters: wave height and wave direction, and then applied a novel data cleaning method, namely Kwon formulas, to remove noise from the data. The results revealed that Kwon formulas are a promising approach to improving data quality. However, it could lead to idealized data used for training, which would lead to poor real-world applicability. In addition, 
they used model test data, a scaled-down representation that is good for predictions in the ship design phase but not in the operational performance context.

In contrast, Zhou et al.~\citep{zhouPredictingShipFuel2023} used in-situ data, which is full-scale data based on the natural environment, e.g., open sea. 
Their model is a 2-stage sequential model based on a fishing vessel. In the first stage, a prediction of \ac{stw} is made based on \ac{sog} and environmental conditions and then using 
\ac{stw} predictions as input for predicting \ac{foc} in the second stage. The final selected model is based on \ac{rf} architecture. Their prediction output is \ac{foc}, which is different than \ac{fc}.
In the sense that \ac{foc} exclusively applies to oil-based fuel types such as \ac{hfo}, which don't take into consideration greener fuel types. This reduces its readiness for future-proof transition to other types of fuels. 
Also, they fused environmental data from \ac{noaa} data store, which lacks global coverage since it's more suited for voyages within \ac{us} coastal waters. In addition, it lacks standardized APIs that streamline data acquisition, making it impractical for real-world applications. 

Agand et al.~\citep{agandFuelConsumptionPrediction2023} investigated different models such as \ac{mlr}, \ac{ann} and ensemble models. They found that ensemble models performed the best, which aligns with Xie et al.~\citep{xieFuelConsumptionPrediction2023} findings. 
Their model is based on a passenger ferry within a 2-year span and in-service data, which is high-frequency, high-precision data that reflects near real-time changes. 
The model predicts hourly \ac{fc}, which makes it, to some extent, suitable for real-time applications such as route optimization. However, they didn't consider localized environmental factors.

Haoqing et al.~\citep{WANG2023104361} attempted to improve \ac{dl} interpretability by building \ac{ann} that integrates physics and domain knowledge in order to have an explainable model that balances 
the trade-off between interpretability and accuracy. They achieved that by using \ac{shap}, a method developed by Scott et al.~\citep{lundberg2017unified} based on game theory. 
The model was trained using publically accessible data over two months from a ferry vessel originating from a study published in Petersen et al.~\citep{petersen2012machine} as cited by Haoqing et al.~\citep{WANG2023104361}. 
The model was trained on very old data, roughly 12 years or more old. This poses a risk of applicability moving forward since the data is no longer relevant, as it won't reflect current shipping practices or take into consideration technologies of new ship builds. 

\subsection{Gaps and Contribution}\label{sec:gaps_and_contribution}

To the best of our knowledge, the most recent and closest work to our study is the work done by Zhang et al.~\citep{zhangDeepLearningMethod2024} and Du et al.~\citep{duDataFusionMachine2022a}. 
Zhang et al.~\citep{zhangDeepLearningMethod2024} developed a model for bulk carriers. Their model is based on \ac{dl} technique, namely \ac{blstm}, and compared it to traditional \ac{ml} methods. 
They found that \ac{blstm} provides better accuracy compared to traditional \ac{dl} and \ac{ml}. However, the difference was not significant enough to justify the 
complexity associated with the usage of such a model in practice. Additionally, they are different to our study that they require high resolution input data and the predictions 
are made per voyage instead of daily predictions. 

Another study closer to ours that was to some extent used as a reference to our study is Du et al.~\citep{duDataFusionMachine2022a}. 
The study's main objective was to fuse environmental conditions to AIS data, and the final model selected is \ac{xgboost}. They are different from our study; however, since they applied it to containerships, our work attempts to apply it in a different context, namely on bulk carriers, which is different due to the data structure and nature of sailing.

Looking at the literature landscape concerning modeling \ac{fc} using \ac{bbm} while fusing external public and reliable meteorological factors, 
To the best of our knowledge, there are not many recent studies that have studied this extensively. Some studies like the work done by Xie et al.~\citep{xieFuelConsumptionPrediction2023} attempted to fuse public environmental parameters but only considering two parameters namely wave heights and direction and therefore no 360° view of the ship envionment was considered.  
This finding is also supported by latest findings by Li et al.~\citep{liDataFusionMachine2022}. 

This raises the overarching question: "Does fusing meteorological factors in ship \ac{fc} modeling gives a better and more accurate model?"

\clearpage